%% notes/topics/90-crosscutting/semisimplicity.tex
%%
%% 用途:
%%   留数・局所モノドロミーの半単純性に関する主張を、3テーマ横断で
%%   照合するための場所。同一の現象が各テーマで異なる地位を持っている。
%%
%% 記入の規則:
%%   - 各テーマでの主張・射程・地位を並記する。
%%   - あるテーマの結果が別のテーマの予想を含意すると見える場合でも、
%%     定義の同値性が確認されるまでは含意を確定させない。
%%     照合に必要な確認事項を明示し、
%%     00-common/literature-status.tex へ登録する。

\section{横断: 半単純性}

\subsection{各テーマでの主張と地位}

\subsection{射程の違い}
%% レベル、局所か大域か、対象の圏。

\subsection{照合に必要な確認事項}

\subsection{照合の結論}
%% 確認が済んでから記入する。未確認のうちは空のままにする。
